\documentclass[11  pt]{article} 
\usepackage[lmargin=1in,rmargin=1.75in,bmargin=1in,tmargin=1in]{geometry}  

\input{preamble}
\usepackage{graphicx}
\usepackage{amsmath}
\usepackage{amssymb}

\begin{document}
\onehalfspacing

\lecture{2: Fast Fourier Transforms}{January 15, 2026}

\paragraph{Course Logistics}

\begin{itemize}
\item Read Chapter 30 of the textbook for this week's material on polynomials and the FFT.
\item Syllabus quiz is due Sat, Jan 17. HW 1 and intro video due Fri, Jan 23.
\end{itemize}

\section{Motivation: Polynomial Multiplication}

We will introduce a powerful application of the divide and conquer paradigm: the Fast Fourier Transform (FFT). We will motivate it through the problem of polynomial multiplication.

\begin{definition} \noindent A \hide{\textbf{polynomial}} in the variable $x$ over an algebraic field is a representation of a function $A(x)$ as a formal sum:
    \begin{equation*}
        A(x) = \sum_{j=0}^{n-1} a_j x^j
    \end{equation*}
    The values $a_0, a_1, \hdots, a_{n-1}$ are the \hide{\textbf{coefficients}} of the polynomial. 
\end{definition}

\vs{.5cm}

\begin{example}\textbf{Polynomial Multiplication} \\
\framebox{
\vbox{
\textbf{Input}: Two polynomials $A(x) = \sum_{j=0}^{n-1} a_j x^j$ and $B(x) = \sum_{j=0}^{n-1} b_j x^j$ of degree bound $n$.
    
\textbf{Output:} A polynomial $C(x) = A(x)B(x)$ of degree bound $2n-1$ such that:
    \begin{equation*}
    C(x) = \sum_{j=0}^{2n-2} c_j x^j \quad \text{where} \quad c_j = \sum_{k=0}^{j} a_k b_{j-k}
    \end{equation*}
}}

\vs{5pt}

An instance of this problem is multiplying $A(x) = 2 + 3x$ and $B(x) = 1 - x$.
\end{example}

\vs{1cm}

\QU What is the asymptotic runtime of multiplying two polynomials of degree bound $n$ using the standard, naive approach?
\begin{enumerate}
    \aitem $O(1)$
    \bitem $O(n)$
    \citem $O(n \log n)$
    \ditem $O(n^2)$
    \eitem Other
\end{enumerate}
\end{Qu}


\newpage
\section{Representations of Polynomials}
To multiply polynomials faster, we can change how we represent them.

\subsection{Coefficient Representation}
We represent $A(x)$ as a vector of coefficients $a = (a_0, a_1, \hdots, a_{n-1})$.
\begin{itemize}
    \item Evaluating $A(x)$ at a given point takes \hide{$O(n)$} time.
    \item Multiplying two polynomials takes \hide{$O(n^2)$} time.
\end{itemize}

\subsection{Point-Value Representation}
A polynomial of degree bound $n$ is uniquely determined by its values at $n$ distinct points. 
We can represent $A(x)$ as a set of $n$ point-value pairs:
\begin{equation*}
    \{ (x_0, y_0), (x_1, y_1), \hdots, (x_{n-1}, y_{n-1}) \}
\end{equation*}
where $y_k = A(x_k)$ for $k = 0, 1, \hdots, n-1$.

\vs{0.5cm}

\textbf{Why use point-value representation?} \\
If we have $A$ and $B$ in point-value form evaluated at the \emph{same} points $x_0, \hdots, x_{2n-1}$, how long does it take to compute $C(x) = A(x)B(x)$ in point-value form? \\

\hide{Since $C(x_k) = A(x_k)B(x_k)$, we just multiply the $y$-values! This takes $O(n)$ time!}

\vs{0.5cm}

\QU How many point-value pairs do we need to compute the polynomial $C(x) = A(x)B(x)$ uniquely, assuming $A$ and $B$ have degree bound $n$?
\begin{enumerate}
    \aitem $n$
    \bitem $2n-1$
    \citem $n^2$
    \ditem $2n$
\end{enumerate}
\end{Qu}

\vs{0.5cm}
\paragraph{The Strategy:}
\begin{itemize}
    \item \textbf{Evaluate}: Convert $A$ and $B$ to point-value form at $2n$ points.
    \item \textbf{Point-wise Multiply}: Multiply the values in \hide{$O(n)$} time.
    \item \textbf{Interpolate}: Convert $C$ back to coefficient form.
\end{itemize}
If we use arbitrary points, Evaluation and Interpolation take \hide{$O(n^2)$} time. We need to choose our evaluation points strategically!

\newpage

\section{Background Math: Complex Roots of Unity}

To evaluate our polynomials efficiently, we will evaluate them at the \textbf{complex roots of unity}.

\begin{definition} \noindent The \textbf{$n$-th complex roots of unity} are the $n$ complex numbers $\omega$ such that:
    \begin{equation*}
        \hide{\omega^n = 1}
    \end{equation*}
\end{definition}

Using Euler's formula, $e^{iu} = \cos(u) + i\sin(u)$, these roots are given by:
\begin{equation*}
    \omega_n^k = e^{2\pi i k / n} \quad \text{for } k = 0, 1, \hdots, n-1
\end{equation*}

\begin{example}\textbf{Elementary Example: $n = 4$} \\
    The 4-th roots of unity are:
    \begin{itemize}
        \item $k=0: \omega_4^0 = $ \hide{$1$}
        \item $k=1: \omega_4^1 = $ \hide{$i$}
        \item $k=2: \omega_4^2 = $ \hide{$-1$}
        \item $k=3: \omega_4^3 = $ \hide{$-i$}
    \end{itemize}
\end{example}

\vs{1cm}

\subsection{Crucial Properties for Divide and Conquer}

\begin{lemma}\textbf{Cancellation Lemma:} For any integers $n \geq 0, k \geq 0, d > 0$:
    \begin{equation*}
        \hide{\omega_{dn}^{dk} = \omega_n^k}
    \end{equation*}
\end{lemma}

\vs{1cm}

\begin{lemma}\textbf{Halving Lemma:} If $n > 0$ is even, then the squares of the $n$ complex $n$-th roots of unity are exactly the $n/2$ complex $(n/2)$-th roots of unity.
    \begin{equation*}
        \hide{(\omega_n^k)^2 = \omega_{n/2}^k}
    \end{equation*}
\end{lemma}

\vs{1cm}
The Halving Lemma implies that if we square our $n$ evaluation points, we are left with only \hide{$n/2$ distinct points}. This is perfect for divide and conquer!

\newpage
\section{Intuition for the Fast Fourier Transform}

We want to evaluate a polynomial $A(x)$ of degree bound $n$ at the $n$-th roots of unity: $1, \omega_n, \omega_n^2, \hdots, \omega_n^{n-1}$. Assume $n$ is a power of 2.

\paragraph{Divide}
We can split $A(x)$ into its even-indexed coefficients and odd-indexed coefficients:
\begin{align*}
    A^{[0]}(x) &= a_0 + a_2 x + a_4 x^2 + \hdots + a_{n-2} x^{n/2 - 1} \\
    A^{[1]}(x) &= a_1 + a_3 x + a_5 x^2 + \hdots + a_{n-1} x^{n/2 - 1}
\end{align*}

We can express $A(x)$ in terms of these new polynomials:
\begin{equation*}
    A(x) = \hide{A^{[0]}(x^2) + x A^{[1]}(x^2)}
\end{equation*}

\paragraph{Conquer}
To evaluate $A(x)$ at the $n$ roots of unity, we need to evaluate $A^{[0]}$ and $A^{[1]}$ at the \emph{squares} of the $n$ roots of unity. \\

By the \textbf{Halving Lemma}, these squares are exactly the \hide{$(n/2)$-th roots of unity}. \\

So we just need to evaluate two polynomials of degree bound $n/2$ at $n/2$ roots of unity. This is two recursive calls of size $n/2$!

\paragraph{Combine}
For $k < n/2$, we compute:
\begin{equation*}
    A(\omega_n^k) = A^{[0]}(\omega_{n/2}^k) + \omega_n^k A^{[1]}(\omega_{n/2}^k)
\end{equation*}

For $k \geq n/2$, we use the fact that $\omega_n^{k + n/2} = -\omega_n^k$:
\begin{equation*}
    A(\omega_n^{k + n/2}) = \hide{A^{[0]}(\omega_{n/2}^k) - \omega_n^k A^{[1]}(\omega_{n/2}^k)}
\end{equation*}

\vs{2cm}

\newpage
\section{The Algorithm}

Here is the pseudocode for the recursive FFT algorithm, which evaluates a polynomial at the $n$-th roots of unity.

\begin{algorithm}
    \caption{Recursive-FFT($a$)}
    \begin{algorithmic}
        \State $n = \textbf{length}(a)$
        \If{$n == 1$}
        \State Return $a$
        \EndIf
        \State $\omega_n = e^{2\pi i / n}$
        \State $\omega = 1$
        \State $a^{[0]} = (a_0, a_2, \hdots, a_{n-2})$
        \State $a^{[1]} = (a_1, a_3, \hdots, a_{n-1})$
        \State $y^{[0]} = \textsc{Recursive-FFT}(a^{[0]})$
        \State $y^{[1]} = \textsc{Recursive-FFT}(a^{[1]})$
        \State Let $y$ be a new array of length $n$
        \For{$k = 0$ to $n/2 - 1$}
            \State $y_k = y^{[0]}_k + \omega \cdot y^{[1]}_k$
            \State $y_{k + n/2} = y^{[0]}_k - \omega \cdot y^{[1]}_k$
            \State $\omega = \omega \cdot \omega_n$
        \EndFor
        \State Return $y$
    \end{algorithmic}
\end{algorithm}

\vs{1cm}

\section{Formal Analysis}

\QU What is the recurrence relation for the runtime $T(n)$ of the \textsc{Recursive-FFT} algorithm?
\begin{enumerate}
    \aitem $T(n) = T(n/2) + \Theta(n)$
    \bitem $T(n) = 2T(n/2) + \Theta(1)$
    \citem $T(n) = 2T(n/2) + \Theta(n)$
    \ditem $T(n) = 4T(n/2) + \Theta(n)$
\end{enumerate}
\end{Qu}

\vs{1cm}

Using the \textbf{Master Theorem} from Lecture 1:
\begin{itemize}
    \item $a = \hide{2}$
    \item $b = \hide{2}$
    \item $f(n) = \hide{\Theta(n)}$
\end{itemize}

\vs{1cm}

Since $f(n) = \Theta(n^{\log_b a})$, we are in Case 2. Thus, the runtime is:
\begin{equation*}
    T(n) = \hide{\Theta(n \log n)}
\end{equation*}

\newpage

\section{Main Takeaway: The Convolution Theorem}

We showed how to evaluate polynomials at roots of unity in $\Theta(n \log n)$ time. To finish polynomial multiplication, we need to interpolate the result back to coefficient form.

\paragraph{Interpolation} 
It turns out that interpolating at the complex roots of unity is mathematically identical to evaluating, but using $\omega_n^{-1}$ instead of $\omega_n$, and dividing the final results by $n$. 

Thus, the \textbf{Inverse FFT} can also be computed in \hide{$\Theta(n \log n)$} time!

\vs{1.5cm}

\paragraph{Fast Polynomial Multiplication}

Given two polynomials $A(x)$ and $B(x)$ of degree bound $n$:
\begin{enumerate}
    \setlength\itemsep{2em}
    \item \textbf{Pad:} Add zeros to $A$ and $B$ so they have degree bound $2n$.
    \item \textbf{Evaluate (FFT):} Compute point-value representations of $A$ and $B$ of length $2n$. \\ \textit{Runtime: \hide{$\Theta(n \log n)$}}
    \item \textbf{Point-wise Multiply:} Compute $y_k = A(\omega_{2n}^k) B(\omega_{2n}^k)$ for all $k$. \\ \textit{Runtime: \hide{$\Theta(n)$}}
    \item \textbf{Interpolate (Inverse FFT):} Convert point-values $y_k$ back to coefficients. \\ \textit{Runtime: \hide{$\Theta(n \log n)$}}
\end{enumerate}

\vs{1cm}

\textbf{Total Runtime for Polynomial Multiplication:} \hide{$\Theta(n \log n)$}. \\

\vs{0.5cm}

By leveraging the \textbf{Divide and Conquer} paradigm and evaluating at mathematically symmetric points, we can bypass the naive $O(n^2)$ limit. FFT is an incredibly important algorithm used heavily in signal processing, image compression, and fast large-integer multiplication!

\end{document}